\documentclass[border=1pt]{standalone}
% Shared preamble for the PixelRoot figures: Times text and math to match
% IEEEtran; every figure is drawn at its printed size (no scaling in the paper).
\usepackage[T1]{fontenc}
\usepackage{mathptmx}
\usepackage{amsmath}
\usepackage{xcolor}
\usepackage{tikz}
\usetikzlibrary{arrows.meta,positioning,calc,fit,backgrounds,decorations.pathreplacing,shapes.geometric}
\definecolor{ink}{HTML}{1B2838}
\definecolor{acc}{HTML}{1F5FBF}
\definecolor{grn}{HTML}{1E7B3A}
\definecolor{red}{HTML}{B3261E}
\definecolor{amb}{HTML}{9A6B00}
\definecolor{cu}{HTML}{B87333}
\tikzset{
  >={Stealth[length=1.6mm,width=1.2mm]},
  every picture/.style={line width=0.6pt, font=\footnotesize, color=ink},
  blk/.style={draw=ink, rounded corners=1.5pt, align=center, inner sep=2.5pt,
    minimum height=9mm, fill=white},
  note/.style={font=\scriptsize, align=left, text=ink!80},
}

\usepackage[american]{circuitikz}
\ctikzset{tripoles/mos style/arrows, transistors/arrow pos=end}
\begin{document}
\begin{circuitikz}[line width=0.6pt]
\ctikzset{bipoles/length=0.9cm}
% photodiode (cathode to the transfer gate)
\draw (0,1.2) node[ground]{} to[photodiode, l_=PD] (0,3.0) -- (0.5,3.0);
% TX: horizontal, gate up
\node[nmos, rotate=-90, anchor=S] (tx) at (0.5,3.0) {};
\node[above=0pt of tx.G] {TX};
\draw (tx.D) -- (2.9,3.0) coordinate (fd);
\fill (fd) circle (1.4pt);
\node[red, anchor=north west, font=\footnotesize\bfseries, inner sep=1pt] at ($(fd)+(0.08,-0.06)$) {FD};
% C_FD (small: high conversion gain)
\draw (fd) to[C] (2.9,1.2) node[ground]{};
\node[anchor=east, align=right] at (2.65,2.1) {$C_\mathrm{FD}$\\[-1pt]{\scriptsize alone: high gain}};
% DCG switch, then RST to VDD
\draw (fd) -- ++(0,0.95) coordinate (dcgs);
\node[nmos, anchor=S] (dcg) at (dcgs) {};
\coordinate (x) at (dcg.D);
\fill (x) circle (1.4pt);
\node[nmos, anchor=S] (rst) at ($(x)+(0,0.5)$) {};
\draw (x) -- (rst.S);
\node[left=0pt of rst.G] {RST};
\draw[acc, <-] (dcg.G) -- ++(-0.35,0) node[left, align=right, font=\scriptsize, text=acc, inner sep=1pt]
  {\textbf{DCG}: keyed\\gain-select bit};
% C_DCG (added: low conversion gain)
\draw (x) -- ++(1.0,0) coordinate (cx);
\draw[amb] (cx) to[C] ++(0,-0.9) node[ground]{};
\node[anchor=west, align=left, text=amb] at ($(cx)+(0.25,-0.45)$) {$C_\mathrm{DCG}$\\[-1pt]{\scriptsize added: low gain}};
% SF + RS readout
\node[nmos, anchor=G] (sf) at (5.9,3.0) {};
\draw (fd) -- (sf.G);
\node[anchor=west] at ($(sf.center)+(0.1,0)$) {SF};
\node[nmos, anchor=D] (rs) at ($(sf.S)+(0,-0.25)$) {};
\draw (sf.S) -- (rs.D);
\node[anchor=west] at ($(rs.center)+(0.1,0)$) {RS};
\draw (rs.S) -- ++(0,-0.3) node[below, font=\scriptsize, inner sep=1pt] {column bus};
% VDD rail
\draw (rst.D) -- ++(0,0.25) coordinate (vddL) -| (sf.D);
\node[above, inner sep=1pt] at ($(vddL)!0.5!(sf.D |- vddL)$) {$V_\mathrm{DD}$};
\end{circuitikz}
\end{document}
